% ============================================================
% qp-fonts.tex —— 字体模块（全品式导学件，供 main.tex \input）
% 职责：西文/中文字体族、黑体楷体族、字重阶梯族（v4.3）、花形斜体特粗族、xeCJK 字符类与断行检查。
% 调用关系：main.tex → qp-fonts → qp-layout → qp-headfoot → qp-titles → qp-blocks。
% 依赖：documentclass [fontset=none]{ctexart}（在 main.tex）。
% v4.3（拍板5 新路线＋执行注意增录1）：
%   正文＝方正书宋 FZShuSong-Z01S（BoldFont=方正黑体 FZHei-B01S）；\heiti＝方正黑体；楷体 KaiTi 维持；
%   数学 Times 系维持。挂载纪律＝xeCJK 一律 Path= 显式路径
%   （C:/Users/28120/AppData/Local/Microsoft/Windows/Fonts/——TinyTeX fontconfig 不扫用户字体目录，
%   fc-list 不可见已实证，禁依赖字体名检索）；方正两款为用户级安装，必须 Path=；
%   KaiTi/Times 属系统目录（fontconfig 可见），维持字体名挂载。
% v4.3 字重阶梯（附则《全品样张细节清单》§三，全品笔画倍数谱→方正黑体分档 FakeBold）：
%   章4.0×＞小节(课时)3.3×＞条目号2.7×＞探究点/例N(◆知识点)2.3×＞题号2.0×＞变式/【】标签1.3×＞正文1×；
%   \heibf 名字保留挂在 2.3× 档（断言⑲沿名）；节标题不加粗（用 \heiti 常规，反向断言）。
% 花形内字（拍板1/总账A）：≈11pt 斜体特粗＝方正黑体＋FakeSlant=0.18（≈10°）＋FakeBold=2.0（fake slant 组合，
%   无现成斜体面，任务书§五-2）。
% ============================================================
\setmainfont{Times New Roman}
\setCJKmainfont[
  Path=C:/Users/28120/AppData/Local/Microsoft/Windows/Fonts/,
  BoldFont=FZHei-B01S.ttf]{FZShuSong-Z01S.ttf}
% v4.4：hei 族补 BoldFont=self——\textbf（题号数字加重，拍板2）在 \heihao 组内对全角「．」走 series=bf，
% 无 BoldFont 声明会回落警告；数字本体走主字体 Times Bold（ BoldFace 自动档），「．」维持 2.0× FakeBold。
\setCJKfamilyfont{hei}[Path=C:/Users/28120/AppData/Local/Microsoft/Windows/Fonts/,BoldFont=FZHei-B01S.ttf]{FZHei-B01S.ttf}
\setCJKfamilyfont{kai}{KaiTi}
% v4.4 拍板4a：仿宋面（说明行→仿宋，全品大题说明行 FZFSK 同位）；系统目录字体按字体名挂载（KaiTi 同例）
\setCJKfamilyfont{fangsong}{FangSong}
% ---- 字重阶梯族（同文件同字库，仅 FakeBold 分档） ----
\setCJKfamilyfont{heizhang}[Path=C:/Users/28120/AppData/Local/Microsoft/Windows/Fonts/,FakeBold=4.0]{FZHei-B01S.ttf}
\setCJKfamilyfont{heijie}[Path=C:/Users/28120/AppData/Local/Microsoft/Windows/Fonts/,FakeBold=3.3]{FZHei-B01S.ttf}
\setCJKfamilyfont{heitiao}[Path=C:/Users/28120/AppData/Local/Microsoft/Windows/Fonts/,FakeBold=2.7]{FZHei-B01S.ttf}
\setCJKfamilyfont{heibf}[Path=C:/Users/28120/AppData/Local/Microsoft/Windows/Fonts/,FakeBold=2.3]{FZHei-B01S.ttf}
\setCJKfamilyfont{heihao}[Path=C:/Users/28120/AppData/Local/Microsoft/Windows/Fonts/,FakeBold=2.0]{FZHei-B01S.ttf}
\setCJKfamilyfont{heibian}[Path=C:/Users/28120/AppData/Local/Microsoft/Windows/Fonts/,FakeBold=1.3]{FZHei-B01S.ttf}
% ---- 花形内字：斜体特粗（fake slant＋FakeBold 组合） ----
\setCJKfamilyfont{huabf}[Path=C:/Users/28120/AppData/Local/Microsoft/Windows/Fonts/,FakeSlant=0.18,FakeBold=2.0]{FZHei-B01S.ttf}
\newcommand{\heiti}{\CJKfamily{hei}}
\newcommand{\heizhang}{\CJKfamily{heizhang}}
\newcommand{\heijie}{\CJKfamily{heijie}}
\newcommand{\heitiao}{\CJKfamily{heitiao}}
\newcommand{\heibf}{\CJKfamily{heibf}}
\newcommand{\heihao}{\CJKfamily{heihao}}
\newcommand{\heibian}{\CJKfamily{heibian}}
\newcommand{\huabf}{\CJKfamily{huabf}}
\newcommand{\kaishu}{\CJKfamily{kai}}
\newcommand{\fangsong}{\CJKfamily{fangsong}}% v4.4 拍板4a：仿宋访问宏（\zhenhead 说明行用；漏挂 0908 首轮补）
\xeCJKDeclareCharClass{CJK}{"2460->"24FF, "2600->"26FF, "2200->"22FF, "25A0->"25FF}
% v4.4（#21 标点收紧）：PunctStyle=plain——全角标点不悬挂不挤压，行尾标点禁出栏、行首禁标点；
%   半角括号/半角标点转换在 postproc_daoxue.py 文本 pass 完成，此处只管断行纪律。
\xeCJKsetup{CheckSingle=true,PunctStyle=plain}
